\chapter{Testing and Evaluation}
\label{ch:evaluation}

\section{Introduction}

This chapter evaluates how well the platform meets its goals. It describes the
evaluation method, then reports the results for retrieval accuracy, answer quality
and honesty, citation quality, understanding of structured documents, and
efficiency. It also summarises the functional testing of the agent and the
knowledge graph. The aim throughout is to judge not only whether the answers are
correct, but also whether the system is honest about what it does not know.

\section{Evaluation Method}

The platform was evaluated in three complementary ways. First, the grounded chat
was assessed on the DOUBLE-BENCH document benchmark \cite{shen2025doublebench},
which poses questions over realistic, multimodal documents and requires one or
several hops of reasoning to answer. This benchmark was chosen because it tests
exactly the setting Synapse targets: retrieval and answering over mixed documents,
including scanned and visually rich pages. Retrieval was measured by whether the
correct source document appeared among the retrieved results, and answers were
graded for correctness.

Second, three purpose-built datasets were created to test the understanding of
structured documents: a set of spreadsheets, a set of Python source files and a
set of C++ source files, each paired with questions whose answers are known. These
let the system's answers be checked exactly against the source, which is harder to
do with free-form prose.

Third, the agent and the knowledge graph were tested functionally, by running
representative tasks end to end and checking that the artefacts and graphs produced
were correct and usable. Because grading free-form answers by hand does not scale,
answer correctness on the benchmark was judged automatically by two independent
language-model judges \cite{zheng2023judge}, and their agreement was recorded so
that the reliability of the grading could itself be assessed.

\section{Retrieval Accuracy}

Retrieval is the foundation of every answer, so it was measured first.
Figure~\ref{fig:eval-retrieval} reports the document-level hit rate, that is, how
often the correct source document was among those retrieved. Retrieval was strong
overall and was essentially perfect on questions that depend on a single document.
It fell on the hardest questions, those that require evidence to be combined across
three documents, which is the expected pattern: the more pieces of evidence a
question needs, the more chances there are to miss one. Retrieval was slightly
lower on scanned documents than on born-digital ones, reflecting the additional
difficulty of reading a scanned page, but the gap was small.

\begin{figure}[H]
  \centering
  \includegraphics[width=0.86\textwidth]{fig_eval_retrieval.png}
  \caption{Document-level retrieval accuracy (hit@5). Retrieval is near-perfect on
  single-document questions and remains strong overall, falling only on the hardest
  questions that require evidence from three documents at once. The small gap
  between born-digital and scanned documents shows that the pipeline reads scanned
  pages effectively. The dashed line marks the overall rate.}
  \label{fig:eval-retrieval}
\end{figure}

\section{Answer Quality and Honesty}

Answer correctness was graded by two independent judges, whose results are shown in
Figure~\ref{fig:eval-answers}. Both judges found that a substantial share of
answers were fully correct, with a further portion partially correct, and the two
judges agreed on three quarters of their verdicts, which gives reasonable
confidence in the grading. The judges differed in strictness, as automatic judges
often do, but they agreed on the overall picture: the system answers a large
fraction of difficult, multi-hop questions correctly while still leaving room to
improve on the hardest cases.

\begin{figure}[H]
  \centering
  \includegraphics[width=0.86\textwidth]{fig_eval_answers.png}
  \caption{Answer correctness as graded by two independent judges. Each bar shows
  the share of answers judged correct, partially correct and incorrect. The two
  judges agreed on about three quarters of their verdicts; they differ in
  strictness but agree that the system answers a large share of difficult,
  multi-hop questions correctly.}
  \label{fig:eval-answers}
\end{figure}

Correctness alone does not capture honesty. A system that answers every question,
right or wrong, is less trustworthy than one that declines when it should. The
overconfidence rate, the share of answers that were given confidently but were
wrong, was low at $0.148$, and applying the abstention threshold designed in
Chapter~4 roughly halved it to $0.074$. In other words, allowing the system to
abstain removed about half of its confident mistakes, which is exactly the
behaviour the design aimed for.

\section{Citation Quality}

Because every claim is meant to carry a citation, the citations themselves were
measured. Table~\ref{tab:citation} reports their precision and recall. Precision
was perfect, meaning that every passage the system cited was genuinely relevant to
the claim it supported, and recall was high, meaning that almost every claim that
needed a citation received one. Reliable citations are what allow a user to trust,
and to check, the answers.

\begin{table}[H]
  \centering
  \caption{Citation quality. Precision measures whether cited passages are
  genuinely relevant; recall measures whether claims that need a citation receive
  one. Perfect precision and high recall mean the citations can be relied upon to
  verify an answer.}
  \label{tab:citation}
  \small
  \begin{tabularx}{0.7\textwidth}{@{}Xcc@{}}
    \toprule
    \textbf{Metric} & \textbf{Precision} & \textbf{Recall} \\
    \midrule
    Citation quality & $1.000$ & $0.963$ \\
    \bottomrule
  \end{tabularx}
\end{table}

\section{Understanding Structured Documents}

The three purpose-built datasets tested whether the system can answer exactly from
structured documents. Figure~\ref{fig:eval-datasets} reports the share of questions
whose answer matched the known value. Accuracy was high on all three: the system
read spreadsheets, Python and C++ files and answered most questions exactly. The
few failures were concentrated on questions that require aggregating across many
rows or that depend on a passage the retrieval step did not surface, which points
to where the most useful improvements lie.

\begin{figure}[H]
  \centering
  \includegraphics[width=0.66\textwidth]{fig_eval_datasets.png}
  \caption{Exact-match accuracy on the purpose-built structured datasets. The
  system answered most questions about spreadsheets, Python files and C++ files
  exactly against the known value. The remaining failures were mainly questions
  needing aggregation across many rows or a passage that was not retrieved, which
  indicates the clearest avenues for improvement.}
  \label{fig:eval-datasets}
\end{figure}

\section{Efficiency}

Response time matters for a usable product, so the latency of the main operations
was measured. Table~\ref{tab:efficiency} reports representative figures. Retrieval
returns its first evidence in about two seconds, so the interface feels responsive.
A full, deep answer that reasons over several passages takes longer, with a typical
time of around half a minute and a long tail on the most demanding questions, which
is the expected cost of the careful, multi-step reasoning that makes the answers
reliable.

\begin{table}[H]
  \centering
  \caption{Representative response times. Retrieval is fast enough for the
  interface to feel immediate, while a full multi-step answer takes longer, with a
  long tail on the most demanding questions. The extra time buys the careful
  retrieval and grading that keep the answers grounded.}
  \label{tab:efficiency}
  \small
  \begin{tabularx}{0.86\textwidth}{@{}Xccc@{}}
    \toprule
    \textbf{Operation} & \textbf{Median} & \textbf{90th pct.} & \textbf{99th pct.} \\
    \midrule
    Evidence retrieval        & $2.0$\,s  & --        & --        \\
    Full multi-step answer    & $32$\,s   & $83$\,s   & $143$\,s  \\
    \bottomrule
  \end{tabularx}
\end{table}

\section{Functional Testing}

The agent and the knowledge graph were tested by running representative tasks from
end to end. Table~\ref{tab:functional} summarises the checks and their outcomes.
The agent produced correct charts and documents, asked for clarification only when
a request was genuinely ambiguous, and kept chart values faithful to the source.
The knowledge graph built successfully from a processed library, drew the entities
and their relationships, and kept each element linked back to its source. All of
the checks passed.

\begin{table}[H]
  \centering
  \caption{Functional testing of the agent and knowledge graph. Each
  representative task was run end to end and checked against the expected
  behaviour. All checks passed, confirming that the artefact and graph features
  behave as designed.}
  \label{tab:functional}
  \small
  \begin{tabularx}{\textwidth}{@{}p{0.07\textwidth}Xc@{}}
    \toprule
    \textbf{No.} & \textbf{Check} & \textbf{Result} \\
    \midrule
    T-1 & A chart request produces a chart whose values match the source file. & Pass \\
    T-2 & A document request produces a readable, downloadable document. & Pass \\
    T-3 & An ambiguous request leads the agent to ask one clarifying question. & Pass \\
    T-4 & A clear request proceeds without an unnecessary question. & Pass \\
    T-5 & A poor chart choice is carried out with a recommendation for a better one. & Pass \\
    T-6 & The knowledge graph builds from a processed library and renders. & Pass \\
    T-7 & Graph nodes and edges link back to their source passages. & Pass \\
    T-8 & Team members see new messages and artefacts without refreshing. & Pass \\
    \bottomrule
  \end{tabularx}
\end{table}

\section{Discussion}

The results show that the design choices paid off. Strong retrieval and reliable
citations mean the answers are grounded and checkable, and the abstention
mechanism measurably reduced confident mistakes, which is the behaviour that makes
the system trustworthy. The weakest area, the hardest multi-hop questions, is the
same area that is hard for all current systems, and the failures on the structured
datasets point clearly at aggregation and retrieval coverage as the next things to
improve. The latency results confirm that the interface stays responsive while the
deeper reasoning runs in the time a user will reasonably wait.

\section{Summary}

This chapter evaluated the platform on a document benchmark, on purpose-built
structured datasets and through functional testing. Retrieval was strong, answers
were correct on a large share of difficult questions, citations were reliable, and
the abstention mechanism roughly halved the rate of confident mistakes. The agent
and knowledge-graph features passed every functional check. The next chapter draws
conclusions and sets out directions for future work.
